% Chapter 1 -- Introduction
% Status: basic draft. Numbers come from tables/numbers.tex. Pending items are marked.
\chapter{Introduction}
\label{ch:introduction}

\section{Background and Motivation}
\label{sec:intro-background}

Recommender systems rank items for a user from that user's past behaviour and from the
behaviour of other users. Collaborative filtering (CF) is the dominant approach: a user and an
item are embedded as vectors and the score of a pair is a function of the two vectors
\cite{rendle2009bpr,he2017ncf}. Training such a model has traditionally required collecting all
interaction histories on a server. Federated learning (FL) removes that requirement: the data
stay on the users' devices, each device trains on its own data, and only model updates are sent
to a server that aggregates them \cite{mcmahan2017fedavg,kairouz2021advances}. Federated
recommendation (FedRec) applies this idea to CF
\cite{ammaduddin2019fcf,chai2021fedmf,zhang2025survey}.

Recommendation is an unusually personal task for FL. In standard federated learning one global
model is expected to serve every client; in recommendation every user is a different task, and
the user vector is by construction specific to one client. The field has therefore converged on
\emph{personalised} federated learning \cite{arivazhagan2019fedper,deng2020apfl,collins2021fedrep},
but it has not converged on \emph{where} the personal part should live. Four placements are
common: no private part at all, with the user vectors averaged like every other parameter; a
private user vector that never leaves the device, as in split learning
\cite{singhal2021fedrecon}; a private scoring function on top of a shared item table, as in
PFedRec \cite{zhang2023pfedrec}; and a blend between a private part and population knowledge,
as in APFL \cite{deng2020apfl} and FedTHE \cite{jiang2023fedthe}. Published comparisons of these
placements are hard to read across papers because the evaluation protocols differ in the
participation fraction, in whether the client fine-tunes the item table before scoring, in the
validation split, and in subtle leaks of the test item into training
\cite{krichene2020sampled}.

The setting that matters in practice is \emph{cross-device}: thousands of clients, one user per
client, a small fraction of them participating in each round, and client state that persists
between participations \cite{kairouz2021advances,hard2018gboard}. Under such a protocol the
private part of a client that has not participated for many rounds becomes stale with respect
to the current global part, and the evaluation itself becomes a source of error: a client that
is scored with the wrong local state silently produces a near-random result. These two
properties, staleness and stateful evaluation, are absent from the full-participation protocols
under which most federated recommendation methods are reported.

\section{Problem Statement}
\label{sec:intro-problem}

This thesis started from a specific method: an adaptive, hierarchically-conditioned
personalisation level $\alpha_u$ that blends each user's vector with a population prototype and
adds a per-user neural scorer, intended to help users with few interactions. Under a cross-device
protocol with $N=\NClients$ clients and $C=\ParticipationC$ participation, the method did not
beat the simpler placements, least of all on sparse users. Rather than discard the work, the
thesis was reorganised around the question the negative result raises:

\begin{quote}
\emph{Where should personalisation live in a cross-device federated recommender, what does each
placement cost in privacy and communication, and how must stateful clients be evaluated so that
the comparison can be trusted?}
\end{quote}

The problem has three parts. First, the four placements must be compared on the same data, the
same protocol and the same per-user scoring so that the differences are attributable to the
placement. Second, the comparison with the strongest published method (PFedRec) must separate the
effect of the protocol, in particular the participation fraction and the client-side post-tuning
of the item table at inference, from the effect of the method, because the published number
(\PublishedPfedrecNdcg{} \ndcg{} on MovieLens 1M) is obtained under a different protocol. Third,
the evaluation of stateful clients must be shown to be valid: every number must come from a
warm, reproducible scoring of all \DsUsers{} users.

\section{Objectives}
\label{sec:intro-objectives}

The general objective is an empirical, statistically supported answer to the placement question
under a locked cross-device protocol, together with the evaluation tooling that makes the answer
reproducible. The specific objectives are:

\begin{enumerate}[label=O\arabic*, leftmargin=2.5em]
  \item to implement four federated CF methods, one per placement, on a common framework with a
        common data split, candidate sampling and scoring routine;
  \item to compare them per user on \ndcg{}, \hr{} and \mrr{}, overall and by user activity
        group, with paired statistical tests;
  \item to measure, rather than assume, the causes of the gap between the PFedRec reproduction
        and its published number;
  \item to analyse the adaptive method as a negative result with measured causes;
  \item to quantify the privacy and communication cost of each placement under an
        honest-but-curious server;
  \item to establish and document an evaluation procedure for stateful clients, including an
        offline re-scoring tool that reproduces the federated evaluation exactly.
\end{enumerate}

\section{Research Questions and Hypotheses}
\label{sec:intro-rq}

\begin{description}[leftmargin=3em, style=nextline]
  \item[RQ1] Which placement of the personal part gives the best ranking quality under the
        cross-device protocol, overall and for sparse users?
  \item[RQ2] Does blending the private part with population knowledge, with a data-driven
        $\alpha_u$, improve over keeping the private part fully local?
  \item[RQ3] How much of the gap between the reproduced and the published PFedRec number is due
        to the scoring protocol (post-tuning and the negative-sampling pool) rather than to the
        participation fraction?
  \item[RQ4] What do the placements cost in uploaded bytes and in what an honest-but-curious
        server can infer about a user?
  \item[RQ5] What evaluation safeguards are necessary for stateful clients, and how large is the
        cost of stale private parts at low participation?
\end{description}

The hypotheses below were formulated before the corresponding experiments were run; H1 is the
original thesis hypothesis and is reported as refuted.

\begin{hypothesis}[H1, original]
The adaptive method beats the three other placements on \ndcg{}, including on sparse users.
\end{hypothesis}
\begin{hypothesis}[H2]
Keeping the user vector private beats averaging it, in an otherwise identical configuration.
\end{hypothesis}
\begin{hypothesis}[H3]
On the PFedRec model, replacing the private head by the population-mean head lowers \ndcg{}, and
no blend $\alpha_u \wu + (1-\alpha_u)\bar{\vect{w}}$ with $\alpha_u<1$ recovers the loss.
\end{hypothesis}
\begin{hypothesis}[H4]
Scoring PFedRec with a client-side post-tuned item table lowers \ndcg{} when the test item can be
drawn as a training negative and raises it when the test item is excluded from the negatives, as
in the official code (registered as predictions A1--A5, Appendix~\ref{app:posttune}).
\end{hypothesis}

\section{Scope and Limitations}
\label{sec:intro-scope}

The study uses one dataset (MovieLens 1M), one random seed for the federated runs, one
cross-device protocol ($N=\NClients$, $C=\ParticipationC$, $R=\NRounds$) and four methods that
differ in more than the placement (loss, embedding dimension, optimiser; Section~\ref{sec:disc-fairness}).
Statistical statements are paired per-user tests on a single run, so they describe variability
across users, not across seeds. The best round of each run is selected on the held-out test item
because the protocol has no validation item; the inflation this causes is reported in
Section~\ref{sec:setup-validity}. The participation fraction is not varied, so the share of the
PFedRec gap that is due to $C$ is bounded only indirectly. \pending{the planned run suite of Section~\ref{sec:res-staleness} adds two seeds for all four
methods; until those runs are in, every statement is single-seed}.

\section{Contributions}
\label{sec:intro-contributions}

\begin{enumerate}[label=C\arabic*, leftmargin=2.5em]
  \item \textbf{A placement study under one cross-device protocol.} Four placements of the
        personal part are compared per user on \DsUsers{} clients with paired tests
        (Chapter~\ref{ch:results}). Keeping the user vector private improves \ndcg{} by
        \GainPersonalizedVsBaselinePct\% over averaging it; a private head on a shared item table
        reaches \NdcgPfedrec; on that model the private head is worth \AbsDeltaHeadGlobalHead{}
        \ndcg{} against the population head and a data-driven blend adds nothing
        ($\alpha^\star=\HeadAlphaStarGlobal$ on a grid of \HeadAlphaGridSize{} values).
        \pending{paired tests across the four methods}.
  \item \textbf{A negative result with measured causes.} The adaptive method is analysed rather
        than discarded: the blending weight pulls taste toward the mean, the per-user network has
        \ParamsPersonalMlp{} parameters for a median of \DsMedianRatings{} interactions, and the
        upload of the user vector weakens the privacy boundary (Sections~\ref{sec:res-adaptive}
        and~\ref{sec:disc-adaptive}).
  \item \textbf{A privacy analysis with numbers.} A threat model for an honest-but-curious
        server, the parameters and bytes uploaded per placement, and a preference-reconstruction
        attack on the uploaded user vector (Section~\ref{sec:res-privacy};
        \pending{byte measurements and attack results}).
  \item \textbf{Valid evaluation of stateful clients.} A cold-state evaluation error that
        reduced \ndcg{} \ColdEvalRatio-fold, the counters and the offline re-scoring tool that
        prevent it, a measured decomposition of the PFedRec reproduction gap, and the cost of
        stale private parts at $C=\ParticipationC$ (Sections~\ref{sec:setup-validity},
        \ref{sec:res-pfedrec} and~\ref{sec:res-staleness}).
\end{enumerate}

\section{Thesis Structure}
\label{sec:intro-structure}

Chapter~\ref{ch:literature} reviews collaborative filtering, federated learning, personalised FL,
federated recommendation and the evaluation pitfalls of the field. Chapter~\ref{ch:method}
formulates the problem, defines the four placements and the four methods, the offline re-scoring
procedure and the threat model. Chapter~\ref{ch:setup} describes the data, the locked protocol,
the implementation, the evaluation safeguards and the statistical analysis.
Chapter~\ref{ch:results} reports the results, Chapter~\ref{ch:discussion} assesses the
hypotheses and the threats to validity, and Chapter~\ref{ch:conclusion} concludes. The
appendices contain the full post-tuning tables, the configurations and the reproducibility
record.
